\section{Related Works}

\parab{Hybrid simulation.} We are not the first to consider integrating an event-driven simulator with a real system. In the 90s, the Wisconsin Wind Tunnel (WWT)~\cite{reinhardt1993wwt, mukherjee2000wwt} explored an execution-driven simulation approach to estimate performance of cache-coherent shared memory systems. WWT also encountered the problem of synchronizing simulated time with direct-execution. The approach adopted in WWT is to simulate every quantum, which is calculated from the minimum inter-core communication latency. In theory, we could also build \sys using this idea by carefully controlling the execution of containers (using interrupts to stop/resume container execution), so that containers' times are synchronized with the simulator. However, this approach would make \sys significantly slower. Another key difference between \sys and execution-driven simulation from the computer architecture community~\cite{martin2005gems} is that \sys does not faithfuly execute data operations. \sys does not track the GPU buffer content and does not execute GPU computation (except profiling runtime) or communication.

Another effort is to enable today's event-driven network simulator with real systems. For instance, ns-3~\cite{ns3} supports a feature called TapBridge, which implements a special Linux network device. This allows Linux applications to run over an ns-3 simulated network. However, this approach introduces inaccuracies in performance estimation. ns-3 does not control the system time of the Linux environment in the same way \sys controls the time of the ML system. As a result, if ns-3 takes 1 second to simulate 10 ms of network activity, the Linux application perceives the network as 100 times slower than it actually is. 

\parab{Predicting CUDA kernel execution time.} Predicting a CUDA kernel's execution time is a standard problem in ML compilers~\cite{chen2018tvm, zheng2020ansor}. The reason is that an ML compiler's goal is to generate the most efficient CUDA kernel for ML operations.
The approaches compilers take are usually ML-based performance modeling, which eliminates the need to test every CUDA kernel's performance on real hardware. For example, TVM \cite{chen2018tvm} trains a gradient boosting decision tree by applying a set of manually-designed features on generated code. However, our problem of predicting CUDA kernel performance differs significantly from that of ML compilers. We focus on a limited set of kernels---those already selected by the ML systems (\eg, from PyTorch, Megatron \cite{shoeybi2019megatron}, DeepSpeed \cite{rasley2020deepspeed} and TorchTitan \cite{liang2025torchtitan}).
So, \sys can use comprehensive runtime profiling of these CUDA kernels instead of other methods to estimate their performance.


\parab{Emulating the control plane and simulating the data plane.} We draw inspirations from CrystalNet \cite{liu2017crystalnet} for network testing. CrystalNet's idea is to emulate a network environment for switch control programs without actually forwarding data. This allows CrystalNet to fully emulate large production networks on only tens of VMs/containers to find network bugs.
Similarly, in \sys, all the ML system code are running as is except the GPU operations and communication are simulated. Although in different contexts, both CrystalNet and \sys rely on the assumption that the control flow does not have data dependency.
